\documentclass[conference]{IEEEtran}
\IEEEoverridecommandlockouts

\usepackage{cite}
\usepackage{amsmath,amssymb,amsfonts}
\usepackage{algorithmic}
\usepackage{graphicx}
\usepackage{textcomp}
\usepackage{xcolor}
\usepackage{enumitem}

\def\BibTeX{{\rm B\kern-.05em{\sc i\kern-.025em b}\kern-.08em
    T\kern-.1667em\lower.7ex\hbox{E}\kern-.125emX}}

\begin{document}

% Working draft for CNSM 2026.
% Editorial direction:
% - focus on the scientific claim: HLD-driven bootstrap of a network from T0
% - keep multi-agent / graph / memory as implementation means, not headline claims
% - integrate related work inside Section II to stay compatible with the 9-page limit

\title{HLD-as-Source-of-Intent: Bootstrapping Intent-Based Networks from High-Level Design}

\author{
\IEEEauthorblockN{Marc Bruy\`ere-Yamada}
\IEEEauthorblockA{
\textit{T\'el\'ecom SudParis} \\
\textit{JFLI} \\
France / Japan \\
email TBD
}
\and
\IEEEauthorblockN{Co-authors TBD}
\IEEEauthorblockA{
\textit{Affiliations TBD} \\
email TBD
}
}

\maketitle

\begin{abstract}
Intent-Based Networking (IBN) is usually studied as an operational layer applied to an already deployed network \cite{}. This draft paper argues that the same intent-centric reasoning can be pushed earlier in the lifecycle, using a High-Level Design (HLD) document as the initial source of intent for bootstrapping a network from design to first deployment. We target small and medium enterprise networks and propose a closed-loop workflow that ingests an HLD, derives implementation artefacts, validates them in a lab or digital-twin environment, and produces a first operational network with human approval at critical steps \cite{}. The prototype relies on an agentic orchestration harness, a structured source of truth, and vendor-aware rendering, but these are implementation choices rather than the paper's primary claim. The main contribution is the scientific positioning of ``HLD-as-intent'' at time zero, its grounding in IBN concepts from RFC~9315 and RFC~9316 \cite{}, and an evaluation methodology across multiple HLD instances compatible with CNSM's focus on automation, intent, and agentic intelligence.
\end{abstract}

\begin{IEEEkeywords}
intent-based networking, high-level design, network automation, closed-loop management, network design, digital twin, agentic systems
\end{IEEEkeywords}

\section{Introduction}

Intent-Based Networking (IBN) is most often introduced once a network already exists: operators define intents for a running infrastructure, and the system translates them into operational actions such as deployment, assurance, and remediation \cite{}. In practice, however, enterprise networks start earlier with human-written design artefacts: requirement documents, architecture decisions, and especially High-Level Designs (HLDs). This paper asks whether that HLD can become the first machine-actionable source of intent.

\begin{itemize}[leftmargin=*,noitemsep,topsep=2pt]
\item \textbf{Problem statement.} There is a gap between the HLD, which captures architecture and rationale, and the deployed network, which is usually produced through manual, vendor-specific implementation steps.
\item \textbf{Scientific angle.} Existing IBN work largely emphasizes operation on an existing network; we instead focus on the \emph{T0} phase, where no network exists yet and the challenge is to bootstrap a first operational state from design-time intent \cite{}.
\item \textbf{Working concept.} We treat this early-stage intent as a form of ``Intent Zero'': intent before operations, before assurance, and before the network has become a managed object.
\item \textbf{Scope choice.} This paper deliberately limits itself to design-to-initial-deployment. A later operational IBN loop may coexist with this approach, but is not the primary object of study here.
\end{itemize}

The paper is built around the following claim: an HLD, provided in a structured and parsable form, can act as the initial source of intent for an end-to-end workflow that produces a first deployable network configuration, validates it in a controlled environment, and prepares a network for subsequent operation and management.

\subsection{Draft Contributions}

\begin{itemize}[leftmargin=*,noitemsep,topsep=2pt]
\item \textbf{C1.} A reframing of IBN from \emph{operations on an existing network} to \emph{bootstrapping a network from an HLD}.
\item \textbf{C2.} A lifecycle view that positions the HLD as the earliest explicit intent artefact, bridging design, dimensioning, implementation planning, and initial deployment.
\item \textbf{C3.} A prototype closed-loop workflow, grounded in RFC terminology, that transforms HLD information into deployment artefacts, validation steps, and as-built feedback \cite{}.
\item \textbf{C4.} An evaluation methodology based on multiple SME HLDs, vendor and constraint profiles, and human-in-the-loop checkpoints.
\end{itemize}

\subsection{Paper Structure}

\begin{itemize}[leftmargin=*,noitemsep,topsep=2pt]
\item Section~II provides background, motivation, and the positioning against related work.
\item Section~III presents the proposition at a high level.
\item Section~IV details the input model, workflow, and implementation choices.
\item Section~V defines the experimental evaluation plan.
\item Section~VI discusses scope boundaries, limitations, and implications.
\item Section~VII concludes.
\end{itemize}

\section{Background and Motivation}

\subsection{Network Lifecycle and the Role of HLD}

\begin{itemize}[leftmargin=*,noitemsep,topsep=2pt]
\item Typical lifecycle: business needs $\rightarrow$ requirements $\rightarrow$ architecture $\rightarrow$ HLD $\rightarrow$ LLD $\rightarrow$ implementation planning $\rightarrow$ deployment $\rightarrow$ validation $\rightarrow$ operations.
\item In current engineering practice, the HLD is often authoritative for humans but inert for machines.
\item The HLD usually contains the architectural intent that matters most at T0:
site structure, service zones, segmentation goals, interconnection principles, resilience expectations, and acceptance criteria.
\item The translation from HLD to device-ready artefacts remains costly, manual, and error-prone in multi-vendor contexts \cite{}.
\end{itemize}

\subsection{IBN Concepts Relevant to This Paper}

\begin{itemize}[leftmargin=*,noitemsep,topsep=2pt]
\item RFC~9315 defines IBN as a closed-loop system that converts high-level intent into network behavior and continuously checks conformance \cite{}.
\item RFC~9315 also distinguishes an inner control loop and an outer human loop; this distinction is useful here, but our paper concentrates on the earliest fulfilment stages rather than steady-state operation \cite{}.
\item RFC~9316 can be used to clarify what kind of intent is involved in this paper:
business intent, service intent, resource intent, and the transformations between them \cite{}.
\item One key question for the paper is therefore not only ``what is intent?'', but also ``at what point in the lifecycle does intent become operationally actionable?''.
\end{itemize}

\subsection{Related Work and Research Gap}

\begin{itemize}[leftmargin=*,noitemsep,topsep=2pt]
\item \textbf{Classical IBN literature and platforms.} Most approaches assume an already deployed network and focus on policy translation, closed-loop assurance, or service steering \cite{}.
\item \textbf{Recent CNSM/NOMS work.} Existing work covers intent translation, traffic steering, NFV or Kubernetes deployment, and KPI assurance with LLMs, but does not clearly address the HLD-to-network bootstrap problem \cite{}.
\item \textbf{MALT-inspired abstraction.} Multi-layer modelling work such as MALT helps justify the need for several abstraction levels across a network lifecycle, but does not provide an open demonstrator for turning HLD content into an initial running network \cite{}.
\item \textbf{Gap.} We did not identify a clear scientific treatment of HLD-as-source-of-intent for initial network realization from a white-field starting point.
\item \textbf{Positioning.} Our paper is therefore closer to \emph{design-time IBN} or \emph{Intent Zero} than to classical operational IBN.
\end{itemize}

\section{Proposition: Overview}

% Figure idea:
% end-to-end view from HLD -> normalized intent -> implementation artefacts
% -> lab/digital twin deployment -> validation -> first operational network.

\subsection{Core Proposition}

\begin{itemize}[leftmargin=*,noitemsep,topsep=2pt]
\item \textbf{Input.} A structured HLD describing the target SME network, possibly complemented by explicit technology constraints.
\item \textbf{Output.} A first validated network realization: topology artefacts, vendor-aware configurations, deployment plan, and as-built feedback.
\item \textbf{Central hypothesis.} A sufficiently structured HLD can serve as the first formalized intent artefact without requiring an operator to re-enter the same intent into a separate controller-specific model.
\item \textbf{Boundary.} The paper does not claim to solve the whole lifecycle; it addresses the transition from design intent to first deployable and validated network state.
\end{itemize}

\subsection{Closed-Loop View}

\begin{itemize}[leftmargin=*,noitemsep,topsep=2pt]
\item Step 1: ingest and parse the HLD into structured design objects.
\item Step 2: normalize those objects into an internal intent representation.
\item Step 3: derive policy, topology, addressing, segmentation, and implementation artefacts.
\item Step 4: validate the result in a lab or digital-twin environment before production push \cite{}.
\item Step 5: generate deployment and validation outputs for human review when changes are high impact.
\item Step 6: feed the observed result back into an as-built view that can later support operations.
\end{itemize}

\subsection{Why a Closed Loop at T0?}

\begin{itemize}[leftmargin=*,noitemsep,topsep=2pt]
\item Because design artefacts and implementation artefacts drift even before the network enters normal operation.
\item Because a design-time loop can catch mismatches early:
missing constraints, unsupported vendor features, invalid topology assumptions, and inconsistent security intent.
\item Because a digital twin or lab validation stage can reduce the risk of turning an HLD directly into production changes \cite{}.
\item Because this first loop can later hand over a cleaner, better documented network to a classical operational IBN loop.
\end{itemize}

\section{Details}

\subsection{Input Model: HLD or Structured Input File}

\begin{itemize}[leftmargin=*,noitemsep,topsep=2pt]
\item Current working assumption: the entry artefact is still called an HLD, but it must contain structured sections that are parsable and reproducible.
\item Candidate sections:
scope, sites, zones, topology intent, addressing, segmentation, external connectivity, security rules, critical services, vendor or platform constraints, validation expectations.
\item Open design choice to discuss in the paper:
whether technology constraints belong inside the HLD or in a companion constraint file.
\item Candidate formalizations to compare in the future:
Markdown with tables, PlantUML, Mermaid, or another machine-readable design notation \cite{}.
\end{itemize}

\subsection{Transformation Workflow}

\begin{itemize}[leftmargin=*,noitemsep,topsep=2pt]
\item \textbf{Design parsing.} Extract structured entities and relations from the HLD.
\item \textbf{Intent normalization.} Convert raw HLD content into a stable internal representation that distinguishes invariant design choices from platform-dependent constraints.
\item \textbf{Policy and topology derivation.} Produce network intents such as segmentation, reachability, redundancy, and service exposure.
\item \textbf{Implementation synthesis.} Render vendor-aware artefacts, low-level configurations, topology descriptors, and deployment plans.
\item \textbf{Validation.} Deploy in Containerlab or an equivalent environment, execute validation checks, and collect discrepancies.
\item \textbf{As-built feedback.} Record the realized state and discrepancies to support traceability and later operation.
\end{itemize}

\subsection{Prototype Realization Choices}

\begin{itemize}[leftmargin=*,noitemsep,topsep=2pt]
\item The current prototype uses an agentic orchestration harness to decompose the workflow into specialized tasks.
\item The present draft intentionally treats this multi-agent decomposition as a \emph{means of realization}, not as the paper's main scientific claim.
\item A graph-backed source of truth helps preserve cross-layer relations between design objects, implementation artefacts, deployment state, and validation results \cite{}.
\item Additional memory layers can be useful for rationale capture, traceability, and human interaction, but should remain secondary in the paper narrative.
\item Vendor-specific rendering and deployment backends are required to demonstrate that the approach is not limited to a single platform \cite{}.
\end{itemize}

\subsection{Human-in-the-Loop and Governance}

\begin{itemize}[leftmargin=*,noitemsep,topsep=2pt]
\item The workflow should remain largely automated for low-risk transformations.
\item A human checkpoint is required for high-impact changes:
new devices, topology modifications, security-sensitive exposure, or unresolved implementation ambiguity.
\item One evaluation dimension is therefore not only technical correctness, but also the \emph{amount and type of human intervention} still required.
\item This is also where closed-vendor documentation, unsupported models, or ambiguous HLD wording become visible.
\end{itemize}

\section{Experimental Evaluation}

% Table idea:
% HLD scenarios, vendor profiles, constraints, expected services, and outcomes.

\subsection{Objectives and Research Questions}

\begin{itemize}[leftmargin=*,noitemsep,topsep=2pt]
\item \textbf{RQ1.} Can a structured HLD provide enough information to generate a first deployable SME network without manual re-entry into another intent model?
\item \textbf{RQ2.} How robust is the workflow across multiple HLDs, topology shapes, and vendor or technology constraints?
\item \textbf{RQ3.} Where is human intervention still required, and is that intervention due to HLD ambiguity, missing constraints, or implementation limitations?
\item \textbf{RQ4.} To what extent are the results dependent on a specific model, toolchain, or vendor backend?
\end{itemize}

\subsection{Experimental Setup}

\begin{itemize}[leftmargin=*,noitemsep,topsep=2pt]
\item Target domain: small and medium enterprise networks.
\item Validation environment: Containerlab or another comparable digital-twin infrastructure \cite{}.
\item Multiple vendor profiles:
for example software routers, Linux-based network functions, and model-driven platforms.
\item Reproducible pipeline execution from the same HLD input to the same expected artefacts and validation checks.
\end{itemize}

\subsection{Scenario Corpus}

\begin{itemize}[leftmargin=*,noitemsep,topsep=2pt]
\item A corpus of roughly ten HLDs with increasing complexity:
single site, multi-site, DMZ exposure, guest access, service segmentation, redundancy, and vendor heterogeneity.
\item Mix of synthetic and, when available, industrially inspired HLDs.
\item Optional stress cases:
underspecified HLD, conflicting constraints, unsupported vendor feature, and closed-vendor documentation dependency.
\end{itemize}

\subsection{Metrics}

\begin{itemize}[leftmargin=*,noitemsep,topsep=2pt]
\item End-to-end success rate per HLD.
\item Number of manual interventions per scenario and their causes.
\item Time to first validated network realization.
\item Coverage of HLD elements successfully translated into deployment artefacts.
\item Validation pass rate:
connectivity, segmentation, reachability, and service tests.
\item Reproducibility across repeated runs.
\item Optional cost indicators:
tokens, latency, and model dependence when an LLM-assisted component is enabled.
\end{itemize}

\subsection{Comparative and Ablation Angles}

\begin{itemize}[leftmargin=*,noitemsep,topsep=2pt]
\item Compare structured-HLD input with less structured variants to quantify the value of formalization.
\item Compare one combined input file against a split HLD-plus-constraints organization.
\item Compare validation-only digital-twin execution against a path that targets a more realistic deployment backend.
\item Ablate optional components such as memory support or different orchestration strategies to show what is essential to the core claim.
\end{itemize}

\section{Discussion}

\subsection{Scientific Positioning}

\begin{itemize}[leftmargin=*,noitemsep,topsep=2pt]
\item The paper should explicitly state that it is \emph{not} a full replacement for operational IBN.
\item Instead, it introduces a preceding phase:
design-time intent realization, or ``Intent Zero'', which prepares a network before classical assurance and operations begin.
\item This clarifies the apparent tension raised by the draft title:
the system is IBN-inspired and RFC-grounded, but it operates earlier in the lifecycle than most existing IBN work.
\end{itemize}

\subsection{Limits and Threats to Validity}

\begin{itemize}[leftmargin=*,noitemsep,topsep=2pt]
\item Results may depend on the structure and quality of the HLD itself.
\item A lab or digital twin remains a proxy for production infrastructure, even if it is useful for pre-deployment validation.
\item Vendor closedness, missing YANG or API models, and incomplete documentation may require extra human intervention.
\item The chosen orchestration architecture may influence engineering efficiency without changing the main scientific claim.
\end{itemize}

\subsection{Implications}

\begin{itemize}[leftmargin=*,noitemsep,topsep=2pt]
\item For IBN research:
intent should be studied not only as an operational abstraction, but also as a design-time artefact.
\item For practitioners:
a disciplined HLD format may reduce the cost of moving from architecture to implementation.
\item For future work:
the same framework could later be extended toward an operational IBN loop, continuous assurance, and design updates after deployment.
\end{itemize}

\section{Conclusion}

\begin{itemize}[leftmargin=*,noitemsep,topsep=2pt]
\item This paper draft positions the HLD as the earliest explicit source of network intent.
\item The main message is that the gap between design and deployment can be framed as an IBN problem, not only as a documentation or automation problem.
\item The proposed workflow aims to transform a structured HLD into a first validated network realization for SME environments.
\item The next steps are to finalize the HLD formalization choice, complete the experimental campaign across multiple HLDs, and sharpen the discussion on coexistence with later operational IBN loops.
\end{itemize}

% TODO later:
% - refine title after deciding whether "Intent Zero" remains in the paper
% - add Figure 1 (positioning in the network lifecycle)
% - add Figure 2 (workflow from HLD to as-built)
% - decide whether a small related-work table is worth the page cost
% - replace all \cite{} placeholders with actual entries
% - re-enable BibTeX once real citation keys are available

\section*{References}
References will be added after the current \texttt{\string\cite\{\}} placeholders are replaced with actual keys.

\end{document}
